\section{Spectral likelihood decoding}
\label{sec:algorithm}

\subsection{Model and retained information}
Let $A$ be the adjacency matrix of an undirected graph without self-loops. Conditional on $z\in[K]^n$, the upper-triangular entries are independent, with
\begin{equation}
 A_{ij}\sim\operatorname{Bernoulli}(P_{z_i z_j}),\quad i<j,
 \qquad A_{ji}=A_{ij},\quad A_{ii}=0.
 \label{eq:sbm}
\end{equation}
We follow Zhou and Li~\citep{zhouli2020}: $P$ is the symmetric block probability matrix, $p_{aj}=P_{a,z_j}$, and $p_{a*}$ is the connection profile of community $a$. Write $Z_{ia}=\one\{z_i=a\}$, $n_a=\sum_iZ_{ia}$, $\pi_{a,n}=n_a/n$, and
\[
 p=\max_{a,b}P_{ab}.
\]
The known community number $K\ge2$ is any fixed integer. Our uniform model and density conditions are, for fixed positive constants $\pi_0,\kappa,\eta$,
\begin{align}
 n_a&\ge\pi_0n,\qquad P_{ab}\ge\kappa p\quad(a,b\in[K]),\qquad
 \sigma_{\min}(P)\ge\kappa p,
 \label{eq:model-conditions}\\
 p&\le1-\eta,\qquad p=\Omega(\log n/n).
 \label{eq:density-condition}
\end{align}
We may take $\kappa\le1$. Neither $P/p$ nor the community proportions need converge. The probability matrix may vary with $n$ within these bounds. In particular, there is no requirement that $(n/\log n)P$ be fixed.

For $z_i=a$, the expected degree is
\begin{equation}
 \E[\deg(i)\mid z]
   =\sum_b(n_b-\one\{a=b\})P_{ab}.
 \label{eq:expected-degree}
\end{equation}
These expectations lie between $\kappa(n-1)p$ and $np$ and may differ across communities. Thus they are uniformly of order $np$, and \eqref{eq:density-condition} implies that all expected degrees are $\Omega(\log n)$. The density condition permits intermediate densities and constant edge probabilities. The probability-comparability and singular-value conditions ensure a uniformly separated, well-conditioned signal; degree alone does not imply identifiability.

Retain the $K$ eigenpairs of the original adjacency matrix with largest absolute eigenvalues:
\begin{equation}
 AU=U\Lambda,\qquad U^\top U=I_K,\qquad
 \widetilde A=U\Lambda U^\top.
 \label{eq:representation}
\end{equation}
Both signs of signal eigenvalues are allowed. Every subsequent operation uses only $(U,\Lambda)$, including initialization, fitting, replacement and classification.

\subsection{Why spectral information can retain a likelihood ratio}
For testing $z_i=a$ against $z_i=b$ with the other labels and $P$ known, put
\begin{equation}
 w_j^{ab}=\log\frac{p_{aj}(1-p_{bj})}{p_{bj}(1-p_{aj})},\qquad
 c_i^{ab}=\sum_{j\ne i}\log\frac{1-p_{aj}}{1-p_{bj}}.
 \label{eq:oracle-contrast}
\end{equation}
The exact incident-edge Bernoulli LLR is $L_i^{ab}=A_iw^{ab}+c_i^{ab}$. Extending $w_i^{ab}$ by the same block formula does not change it because $A_{ii}=0$. The contrast vector is block constant.

\begin{proposition}[Preservation of likelihood contrasts]
\label{prop:exact-preservation}
With $\Pi=UU^\top$, the identity
\begin{equation}
 \widetilde A=A\Pi,\qquad
 \widetilde A_iw-A_iw=A_i(\Pi-\mathrm{Id})w
 \label{eq:projection-identity}
\end{equation}
holds for every vector $w$. If $w\in\operatorname{col}(U)$, the corresponding linear statistic is preserved exactly. If $P$ is nonsingular, then $\operatorname{col}(ZPZ^\top)=\operatorname{col}(Z)$, which contains every $w^{ab}$.
\end{proposition}

The first identity follows from the eigenvector equation. For the second assertion, use the factorization
$ZPZ^\top=(ZN^{-1/2})(N^{1/2}PN^{1/2})(ZN^{-1/2})^\top$, where $N=Z^\top Z$. Thus the spectral decomposition retains directions on which the likelihood acts; it need not itself optimize a likelihood. The actual mean is
\begin{equation}
 \E[A\mid z]=ZPZ^\top-\operatorname{diag}(P_{z_i z_i}).
 \label{eq:mean}
\end{equation}
The statistical analysis below accounts for this diagonal correction and for the difference between the empirical and population subspaces.

\subsection{Observed profiles and the complete Bernoulli score}
To define logarithms of reconstructed probabilities, set
\begin{equation}
 s=\max\{1,\|\Lambda\|_{\mathrm{op}}\},\qquad
 \varepsilon_n=c_0s/n,\qquad
 q_{ij}=\operatorname{clip}(\widetilde A_{ij},\varepsilon_n,1-\varepsilon_n),
 \label{eq:profiles}
\end{equation}
where $0<c_0<\min\{\kappa,\eta,1\}/8$ is a fixed algorithm parameter and $\operatorname{clip}(x,l,u)=\min\{u,\max\{l,x\}\}$. The admissibility of $c_0$ is an explicit model condition on this chosen logarithm floor. The scale $s$ is calculated from the retained eigenvalues. This operation occurs after eigendecomposition. The reconstructed diagonal is retained. Since $s\le n$ for an adjacency matrix and $c_0<1/8$, the interval is nonempty even on exceptional graphs.

For $\widehat p\in(0,1)^n$, the working score is
\begin{equation}
 \ell_i(\widehat p)
 =\sum_{j=1}^n\{q_{ij}\log\widehat p_j
           +(1-q_{ij})\log(1-\widehat p_j)\}.
 \label{eq:working-score}
\end{equation}
The second term retains the information in absent edges at every density. The $q_i$ are dependent fractional profiles, so this is a working fitting criterion. Its relation to the exact incident-edge likelihood is established by the probability bounds in Section~\ref{sec:main-results}.
For fitted profiles $\widehat p_1,\ldots,\widehat p_K$ and weights $\widehat\pi\in\Delta_K$, fit
\begin{equation}
 \mathcal L(\widehat\pi,\widehat p)
 =\sum_i\log\sum_a\widehat\pi_a e^{\ell_i(\widehat p_a)},\qquad
 \widehat z_i=\argmax_a\{\ell_i(\widehat p_a)+\log\widehat\pi_a\}.
 \label{eq:mixture-and-decode}
\end{equation}
The output error is the permutation-invariant fraction
$\Mis(\widehat z,z)=\min_{\sigma\in\mathfrak S_K}n^{-1}\sum_i\one\{\sigma(\widehat z_i)\ne z_i\}$.

\subsection{Growing initialization and exact updates}
\label{sec:growing-branch}
Start one profile at the global mean $n^{-1}\sum_iq_i$. At stage $t=2,\ldots,K$, if the existing profiles are $\widehat p_1,\ldots,\widehat p_{t-1}$, compute
\begin{equation}
 h_i=\max_{a<t}\ell_i(\widehat p_a),\qquad
 \Delta(c)=\sum_i(\ell_i(q_c)-h_i)_+.
 \label{eq:growing-gain}
\end{equation}
Append a node profile maximizing this all-node gain, excluding previously added node IDs, initialize the $t$ weights uniformly, and jointly fit all $t$ profiles. No observations are removed. The fitted profiles supply the next stage. This is the existing global-mean growing construction with the complete score \eqref{eq:working-score}.

At any component count, the exact E/M updates are
\begin{align}
 R_{ia}&=\frac{\widehat\pi_a e^{\ell_i(\widehat p_a)}}
                  {\sum_b\widehat\pi_b e^{\ell_i(\widehat p_b)}},
 &N_a&=\sum_iR_{ia},\label{eq:responsibilities}\\
 \widehat p_{aj}^{\mathrm{new}}&=\frac{\sum_iR_{ia}q_{ij}}{N_a},
 &\widehat\pi_a^{\mathrm{new}}&=N_a/n.
 \label{eq:mstep}
\end{align}
A component of zero responsibility mass retains its previous profile. Each fit performs at least one M-step. Further exact updates may use any prescribed finite iteration bound and any stopping rule within that bound. All means remain convex combinations of the same $q$ rows, and each E/M pair increases the same working objective. The update formula is unchanged by including the nonedge term: maximizing a weighted Bernoulli cross-entropy still gives the weighted mean.

\subsection{Repeated likelihood replacement for every fixed \texorpdfstring{$K$}{K}}
\label{sec:replacement}
The general-$K$ guarantee uses a deterministic repetition of the existing single-component replacement operation. Set $T_n=\lceil\log n\rceil$. In each round, refer all trials to the same current fit. For each component $r\in[K]$, delete its profile and compute
\begin{equation}
 h_i^{(-r)}=\max_{a\ne r}\ell_i(\widehat p_a),\qquad
 c_r\in\argmax_{c\in[n]}\sum_i(\ell_i(q_c)-h_i^{(-r)})_+.
 \label{eq:replacement-gain}
\end{equation}
Replace profile $r$ by $q_{c_r}$, initialize the trial's weights uniformly, and run a finite exact EM fit. The candidate scan here includes every node ID, including previously used IDs. Compare the $K$ trial endpoints and the unchanged current fit using $\mathcal L$, and keep the largest value. Ties are deterministic, with the current fit retained on an objective tie. The algorithm may stop after a full round with no improving trial; the contraction below certifies sufficiently small loss at that point.

After the replacement rounds, compute the selected fit's responsibilities and perform one mandatory profile and weight M-step before decoding. This last step converts the weak responsibilities certified by the objective into the required local likelihood accuracy. Additional exact EM updates after this step are allowed.

\begin{center}
\fbox{\begin{minipage}{0.93\linewidth}
\textbf{Algorithm 1: Spectral likelihood with repeated parameter replacement}
\begin{enumerate}
\item Form $q$ from the retained eigenpairs by \eqref{eq:profiles}.
\item Obtain $K$ fitted profiles by global-mean growing and joint EM.
\item For $T_n=\lceil\log n\rceil$ rounds, test all $K$ deletion positions; choose each replacement node by \eqref{eq:replacement-gain}, jointly fit it from uniform weights, and retain the greatest working likelihood across the trials and the old fit.
\item Perform one final E-step and profile/weight M-step, then classify by \eqref{eq:mixture-and-decode}.
\end{enumerate}
\end{minipage}}
\end{center}

This is an explicit change from a single replacement pass. The proof gives a loss contraction for all $K$ fitted profiles simultaneously and is valid for every fixed $K$, including unequal communities and either sign of signal eigenvalues. It requires no random restart, no Euclidean initializer, and no assumption that growing alone reaches the correct basin. The unchanged growing-only path is a separate algorithm; its unrestricted general-$K$ convergence is not claimed here.

\subsection{Computational cost}
Forming and storing $q$ requires $O(n^2)$ storage. The complete node-candidate score matrix can be cached using $O(n^3)$ arithmetic and $O(n^2)$ storage. With current component scores cached, all $K$ deletion scans in a round cost $O(Kn^2+K^2n)$: obtain the minimum and second minimum divergences of each row to handle every deletion. If each trial uses at most $J$ EM updates, its fitting cost is $O(JKn^2)$. The $O(\log n)$ rounds therefore cost $O(JK^2n^2\log n)$ in addition to the candidate cache and growing fits. The direct implementation is polynomial in $n$ for fixed $K$ and polynomial $J$. It does not yet provide a subquadratic storage claim.
